\title{Controlling Text-to-Image Diffusion by Orthogonal Finetuning}

\begin{abstract}
Recent advances in large-scale text-to-image diffusion models have enabled highly realistic image synthesis from textual input. However, steering these robust models to perform various downstream applications remains a significant open research area. In response, we propose a systematic finetuning technique termed Orthogonal Finetuning~(OFT) designed for adapting text-to-image diffusion models for diverse tasks. In contrast to prior approaches, OFT mathematically guarantees the preservation of hyperspherical energy—a metric quantifying inter-neuron relationships on the unit hypersphere. Our analysis indicates that maintaining this property is vital to retaining the semantic image generation capacities of diffusion models. Additionally, to enhance the robustness of finetuning, we introduce Constrained Orthogonal Finetuning (COFT), which enforces an added radius constraint on the hypersphere. We focus our study on two primary finetuning scenarios: subject-driven generation, aimed at synthesizing images of a given subject (from limited examples) in novel settings based on text prompts; and controllable generation, which requires the model to accommodate supplementary control signals. Experimental evaluation demonstrates that our OFT framework exceeds current techniques with respect to both image generation quality and speed of convergence.
\end{abstract}

\section{Introduction}

State-of-the-art text-to-image diffusion models~\cite{saharia2022photorealistic,ramesh2022hierarchical,rombach2022high} now deliver outstanding performance in synthesizing high-quality images under textual direction. Nevertheless, reliance on text prompts alone may lead to ambiguities and falls short in enabling precise and detailed manipulations of the generated outputs. To overcome these limitations, this work investigates two specific text-to-image generation challenges:

\begin{itemize}[leftmargin=*,nosep]

    \item \textbf{Subject-driven generation}~\cite{ruiz2023dreambooth}: When only a handful of images for a specific subject are available, this task focuses on producing new images of that subject within varying contexts as instructed by textual input.
    \item \textbf{Controllable generation}~\cite{zhang2023adding,mou2023t2i}: Here, the model is provided with an auxiliary control cue (\emph{e.g.}, edge maps, segmentation labels) alongside a text prompt, and is required to synthesize images in accordance with both inputs.
\end{itemize}

Both of these challenges fundamentally address the problem of adapting text-to-image diffusion models through finetuning, while maintaining the generative capabilities acquired during pretraining. An effective finetuning approach should ideally fulfill two main criteria: (1) \emph{training efficiency}, characterized by a reduced number of tunable parameters and fewer training epochs; and (2) \emph{generalizability preservation}, ensuring that the model retains its ability to generate diverse, high-quality outputs. Standard finetuning practices generally involve either minimally adjusting neuron weights using a low learning rate (\emph{e.g.}, \cite{ruiz2023dreambooth}) or incorporating a small module with re-parameterized weights (\emph{e.g.}, \cite{hulora2022,zhang2023adding}). While straightforward, these finetuning strategies do not reliably maintain the generative performance attained during pretraining. Additionally, there is a lack of a systematic framework for designing optimal finetuning schemes or selecting suitable hyperparameters, such as the number of epochs. One significant obstacle is the absence of a robust metric to assess how well the original generative ability is conserved post-finetuning. Prevailing approaches often presume that minimizing the Euclidean distance between finetuned and pretrained models correlates with better retention of the original performance, which leads to the common practice of applying very low learning rates. Although this assumption is sometimes valid, we posit that solely relying on Euclidean distance falls short in reflecting the full extent of semantic retention. Thus, a more structurally grounded metric to differentiate finetuned models from their pretrained counterparts would offer substantial advantages in preserving pretrained performance and enhancing the robustness of the finetuning process.

Drawing inspiration from empirical studies showing that hyperspherical similarity effectively captures semantic relationships~\cite{liu2017deep,liu2018decoupled,chen2020angular}, our approach employs hyperspherical energy~\cite{liu2018learning} to represent pairwise relationships among neurons within a layer. Specifically, hyperspherical energy quantifies the total hyperspherical similarity (such as cosine similarity) across all neuron pairs in a layer, thereby reflecting the degree of distributional uniformity on the unit hypersphere~\cite{liu2021learning}. We posit that optimal finetuned models should exhibit minimal deviation in hyperspherical energy relative to their pretrained counterparts. A straightforward strategy might involve enforcing a regularization term to keep hyperspherical energy constant throughout finetuning; however, this does not ensure effective minimization of energy differences. Instead, we utilize a key invariance property: applying a uniform orthogonal transformation across all neurons maintains pairwise hyperspherical similarity. Leveraging this characteristic, we introduce Orthogonal Finetuning~(OFT), an approach for adapting large text-to-image diffusion models to downstream tasks while preserving hyperspherical energy. The essence of OFT is to learn a shared orthogonal transformation per layer, thereby maintaining all pairwise angles among neurons. Conceptually, OFT can be interpreted as a rotation or adjustment of the underlying coordinate system within a layer. Further considering that maintaining a smaller Euclidean distance between the finetuned and pretrained model states helps to better retain performance acquired during pretraining, we develop a variant called Constrained Orthogonal Finetuning~(COFT). COFT restricts the updated model parameters to lie on a hypersphere of fixed radius centered at the pretrained neurons.

The rationale behind employing orthogonal transformation for neuron finetuning can be partly traced to insights from the 2D Fourier transform, which allows an image to be represented through its magnitude and phase spectra. The phase spectrum, encoding the angular relationship between the input and the basis, largely captures the image's semantic information. Notably, reconstructing an image with its original phase spectrum paired with a random magnitude spectrum retains its core semantics. This indicates that modifying neuron orientations is essential for inducing semantic changes in generated images—a primary objective in both subject-driven and controllable generative tasks. Nevertheless, unconstrained modification of neuron orientations introduces excessive flexibility, potentially impairing the performance attained through pretraining. To limit this flexibility, we advocate for maintaining the angles between pairs of neurons, premised on the conjecture that these inter-neuron angles are integral to the neural network’s knowledge representation. Accordingly, we propose to optimize a layer-wise, shared orthogonal transformation for neurons within each layer, ensuring that the hyperspherical energy remains constant.

Our methodology also takes cues from orthogonal over-parameterized training~\cite{liu2021orthogonal}, where classification neural networks are constructed by applying orthogonal transformations to randomly initialized models. The motivation lies in the fact that randomly initialized neural networks are expected to possess low hyperspherical energy, and the training approach of \cite{liu2021orthogonal} aims to retain this property, as lower hyperspherical energy correlates with improved classification generalization~\cite{liu2018learning,lin2020regularizing}. The findings in \cite{liu2021orthogonal} demonstrate that orthogonal transformations offer ample expressiveness for training classification models that generalize well. In contrast, our focus is on finetuning text-to-image diffusion models to achieve both enhanced controllability and superior generative capacity in downstream applications. There are two major distinctions between OFT and the method in \cite{liu2021orthogonal}. First, whereas \cite{liu2021orthogonal} targets a reduction in hyperspherical energy, OFT is designed to maintain the hyperspherical energy from the pretrained model, thus preserving its underlying structure during finetuning—a crucial aspect, since lowering the hyperspherical energy for diffusion models might undermine their semantic representations. Second, OFT is formulated to minimize the discrepancy from the original pretrained model, introducing an additional constraint that is absent in the method of \cite{liu2021orthogonal}. Successfully finetuning requires balancing flexibility with model stability, and we contend that the OFT framework achieves an effective equilibrium between these competing demands. The following summarizes our principal contributions:

\begin{itemize}[leftmargin=*,nosep]

    \item We introduce Orthogonal Finetuning, a new finetuning technique aimed at enhancing the controllability of text-to-image diffusion models. For greater stability, we additionally develop a constrained version that restricts angular divergence from the original pretrained parameters.
    \item In contrast to existing approaches, OFT fine-tunes model weights while mathematically guaranteeing conservation of hyperspherical energy, a quantity we find empirically to correlate closely with maintaining the semantic generative capabilities of the original model.
    \item We evaluate OFT on two specific tasks: subject-driven generation and controllable generation. Through extensive experimentation, we show that our method substantially outperforms previous techniques regarding image quality, convergence rate, and robustness during finetuning. Furthermore, OFT demonstrates superior sample efficiency, reaching effective convergence with fewer images and epochs required for training.
    \item For the controllable generation scenario, we present a novel control consistency metric designed to measure controllability. The principle of the metric involves inferring the control signal from the output image and then comparing this inferred signal to the original. This metric provides further evidence of OFT’s high level of controllability.
\end{itemize}

\section{Related Work}

\textbf{Text-to-image diffusion models}. Recent years have witnessed remarkable advancements~\cite{saharia2022photorealistic,ramesh2022hierarchical,rombach2022high,gu2022vector,nichol2022glide} in text-to-image synthesis, primarily driven by the progress in diffusion-based generative frameworks~\cite{ho2020denoising,song2021score,song2021denoising,dhariwal2021diffusion} and multi-modal representation learning~\cite{su2020vlbert,lu2019vilbert,radford2021learning,wang2021simvlm,li2022blip,li2023blip,alayrac2022flamingo,schuhmann2022laion}. GLIDE~\cite{nichol2022glide} and Imagen~\cite{saharia2022photorealistic} focus on pixel-space diffusion. In particular, GLIDE utilizes paired text-image inputs to jointly train a text encoder and diffusion prior, whereas Imagen employs a frozen pretrained text encoder throughout training. Latent-space diffusion models are exemplified by Stable Diffusion~\cite{rombach2022high} and DALL-E2~\cite{ramesh2022hierarchical}; Stable Diffusion leverages VQ-GAN~\cite{esser2021taming} to construct its latent codebook, while DALL-E2 adopts the CLIP~\cite{radford2021learning} joint embedding space for both image and text representations. Beyond diffusion models, generative adversarial networks~\cite{reed2016generative,zhang2017stackgan,xu2018attngan,li2019controllable} and autoregressive architectures~\cite{ramesh2021zero,ding2021cogview,wu2022nuwa,yuscaling} have also contributed to this area. OFT is a universally applicable finetuning methodology and is compatible with any text-to-image diffusion model.

\textbf{Subject-driven generation}. To avoid alterations of the subject, methods such as \cite{avrahami2022blended,nichol2022glide} incorporate user-provided masks as additional conditioning information. Approaches based on inversion~\cite{choi2021ilvr,dhariwal2021diffusion,ramesh2022hierarchical,gal2022image} modify contextual content without affecting the depicted subject. The technique in \cite{hertz2022prompt} supports local and global modifications independently of masks. Nonetheless, these strategies often fall short of fully retaining detailed, identity-specific characteristics. Pivotal Tuning~\cite{roich2022pivotal} proposes local finetuning around an inverted latent code with supplementary regularization to safeguard identity preservation. Analogously, \cite{nitzan2022mystyle} learns a bespoke facial prior using a personal face image dataset. The framework in \cite{casanova2021instance} allows for instance variation generation, though often at the cost of instance-specific detail retention. DreamBooth~\cite{ruiz2023dreambooth} fine-tunes diffusion models with a small number of subject images by introducing a custom token, leveraging both reconstruction and prior preservation losses tied to the subject’s class. While OFT adopts the operational protocol of DreamBooth, it replaces the conventional finetuning—typically with small learning rates—with finetuning via orthogonal transformations for improved performance.

\textbf{Controllable generation}. Image-to-image translation falls under the broader category of controllable generation. Traditionally, most techniques in this area have utilized conditional generative adversarial networks~\cite{isola2017image,zhu2017toward,wang2018high,park2019semantic,choi2018stargan}. In addition, diffusion models have been applied to similar translation tasks~\cite{saharia2022palette,voynov2022sketch,wang2022pretraining}. Recently, ControlNet~\cite{zhang2023adding} demonstrated an effective approach for guiding pretrained diffusion models by fine-tuning and conditioning them on extra control inputs, leading to strong performance in controllable image synthesis. Similarly, T2I-Adapter~\cite{mou2023t2i} adapts diffusion models by fine-tuning for enhanced controllability over generated images. In alignment with the protocol set by \cite{zhang2023adding,mou2023t2i}, we employ OFT to fine-tune pretrained diffusion models, resulting in superior controllability with reduced demands for training samples and a smaller number of parameters updated during fine-tuning. Notably, OFT introduces no additional computational cost during inference.

\textbf{Model finetuning}. The practice of fine-tuning large-scale pretrained models for downstream applications has become highly prevalent~\cite{devlin2018bert,brown2020language,he2022masked}. Within this paradigm, various adaptation strategies (\emph{e.g.}, \cite{hulora2022,houlsby2019parameter,pfeiffer2020adapterfusion}) have been thoroughly explored, particularly in the context of natural language processing. Among these, LoRA~\cite{hulora2022} stands out as highly related to OFT, as it posits a low-rank configuration for additive updates to weights during the fine-tuning process. Distinct from this, OFT applies a shared orthogonal transformation at the layer level to adjust neuron weights multiplicatively, which analytically maintains the relative angles between neurons in a layer, contributing to improved training stability.

\section{Orthogonal Finetuning}

\subsection{Why Does Orthogonal Transformation Make Sense?}

We begin by exploring the rationale for employing orthogonal transformations when fine-tuning text-to-image diffusion models. Specifically, this consideration can be divided into two aspects: (1) the motivation for adjusting the angles (directions) of neurons during fine-tuning, and (2) the justification for utilizing orthogonal transformations to carry out such angular modifications.

In addressing the first question, we draw from prior empirical findings~\cite{liu2018decoupled,chen2020angular}, which suggest that angular differences between features effectively capture the underlying semantic disparity. SphereNet~\cite{liu2017deep} provides evidence that constraining all neuron activations to reside on a unit hypersphere during neural network training not only preserves expressive power but also enhances generalization capability, indicating that the directional component of neuron activations encodes the most salient data characteristics. To more concretely illustrate the relevance of neuron direction, we design a simple experiment as presented in Figure~\ref{fig:toy_ae}. Here, a typical convolutional autoencoder is trained on a dataset of flower images, where the usual inner product is used in the feature map computation during training (with $z$ representing the output of the convolution kernel $\bm{w}$ over an input patch $\bm{x}$). For evaluation, we test three distinct feature map constructions: (a) the conventional inner product from training, (b) isolating only magnitude information, and (c) preserving only angular information. The results, visualized in Figure~\ref{fig:toy_ae}, reveal that the angular aspect of the neurons alone suffices to almost perfectly reconstruct the original images, whereas retaining only the magnitude provides no useful signal. Notably, cosine-based activations are not used in training—the network is trained solely with the standard inner product. This observation strongly suggests that it is predominantly the directional attributes of neuron activations that encode semantic content. Consequently, adapting neuron directions is likely a highly effective approach for modifying the semantic representation within images.

For the second question, the most direct strategy to adjust neuron directions is through collective orthogonal transformations, such as global rotations or reflections applied across all neurons within a particular layer. Although alternative angle-based manipulations—like varying rotation angles for individual neurons—could offer additional flexibility, orthogonal transformations tend to strike an optimal balance between adaptability and structural consistency. Supporting this, \cite{liu2021orthogonal} demonstrates the sufficiency of orthogonal operations in facilitating neural network learning. To empirically corroborate the regularizing effect of orthogonality, we conduct an experiment modeled after the ControlNet~\cite{zhang2023adding} framework, and the findings are illustrated in Figure~\ref{fig:ortho}. Results indicate that our standard Orthogonal Feature Transformation (OFT) consistently achieves stable and precise control after 20 epochs of training, whereas the OFT variant lacking the orthogonality constraint fails to generate plausible outputs or provide meaningful control. These results underscore the crucial role of orthogonality in the effectiveness of OFT.

\subsection{General Framework}

Typically, standard finetuning is carried out by applying gradient descent with a relatively low learning rate to update the model parameters, or sometimes only select layers. The use of a small learning rate implicitly acts as a regularizer, favoring minimal changes from the original pretrained weights; thus, conventional finetuning can be viewed as approximately optimizing the objective $\thickmuskip=2mu \medmuskip=2mu \| \bm{M} - \bm{M}^0\|$, where $\bm{M}$ and $\bm{M}^0$ represent the finetuned and pretrained model weights, respectively. While this approach aligns with intuitive expectations, it might still afford excessive flexibility when applied to finetuning large-scale models. To mitigate this, LoRA augments the finetuning process with an explicit low-rank constraint on the weight modifications, that is, enforcing $\thickmuskip=2mu \medmuskip=2mu \text{rank}(\bm{M} - \bm{M}^0)=r'$ for some small $r'$. In contrast, OFT imposes a condition on the similarity among pairs of neurons, formalized as $\thickmuskip=2mu \medmuskip=2mu \| \text{HE}(\bm{M})-\text{HE}(\bm{M}^0) \|=0$, where $\text{HE}(\cdot)$ stands for the hyperspherical energy of a weight matrix. To illustrate this idea, suppose we have a fully connected layer described by the weight matrix $\thickmuskip=2mu \medmuskip=2mu \bm{W}=\{\bm{w}_1,\cdots,\bm{w}_n\}\in\mathbb{R}^{d\times n}$, with individual neuron vectors $\bm{w}_i\in\mathbb{R}^{d}$, and let $\bm{W}^0$ be the corresponding pretrained weights. The layer output $\thickmuskip=2mu \medmuskip=2mu \bm{z}\in\mathbb{R}^n$ is then given by $\thickmuskip=2mu \medmuskip=2mu \bm{z}=\bm{W}^\top\bm{x}$ for an input $\thickmuskip=2mu \medmuskip=2mu \bm{x}\in\mathbb{R}^d$. Within this framework, OFT can be regarded as the minimization of the hyperspherical energy discrepancy between the finetuned and pretrained networks:

\begin{equation}\label{eq:oft_principle}
\footnotesize
\min_{\bm{W}} \left\| \text{HE}(\bm{W})-\text{HE}(\bm{W}^0)\right\|~~\Leftrightarrow~~\min_{\bm{W}} \bigg{\|} \sum_{i\neq j} \|\hat{\bm{w}}_i-\hat{\bm{w}}_j\|^{-1}-\sum_{i\neq j} \|\hat{\bm{w}}^0_i-\hat{\bm{w}}^0_j\|^{-1}\bigg{\|}
\end{equation}

where each normalized neuron is defined as $\thickmuskip=2mu \medmuskip=2mu \hat{\bm{w}}_i:=\bm{w}_i/\|\bm{w}_i\|$, and the hyperspherical energy for a layer $\bm{W}$ is given by $\thickmuskip=2mu \medmuskip=2mu \text{HE}(\bm{W}):= \sum_{i\neq j} \|\hat{\bm{w}}_i-\hat{\bm{w}}_j\|^{-1}$. It is evident that the optimal value achievable in Eq.~\eqref{eq:oft_principle} is zero, which can be realized whenever the finetuned and original weight matrices differ solely by a rotation or reflection transformation; that is, $\thickmuskip=2mu \medmuskip=2mu \bm{W}=\bm{R}\bm{W}^0$ for some orthogonal matrix $\thickmuskip=2mu \medmuskip=2mu \bm{R}\in\mathbb{R}^{d\times d}$ (where the determinant of $\bm{R}$, $1$ or $-1$, distinguishes between rotation and reflection). This embodies the central principle behind OFT: neural network fin

\begin{equation}\label{eq:oft_general}
    \footnotesize
    \bm{z} = \bm{W}^\top \bm{x} = (\bm{R} \cdot \bm{W}^0)^\top \bm{x}, ~~~ \text{s.t.}~ \bm{R}^\top \bm{R} = \bm{R} \bm{R}^\top = \bm{I}
\end{equation}

Here, $\bm{W}$ corresponds to the weight matrix after OFT-based finetuning, while $\bm{I}$ stands for the identity matrix. An overview of the OFT approach can be found in Figure~\ref{fig:block}. Following a rationale akin to LoRA's zero initialization, it is essential that OFT begins finetuning precisely from the original pretrained parameters $\bm{W}^0$. This is accomplished by setting the initial orthogonal matrix $\bm{R}$ to the identity matrix, thereby ensuring the finetuned parameters match the pretrained ones at the outset. Preserving the orthogonality constraint on $\bm{R}$ can be effectively managed using differentiable orthogonalization techniques as outlined in \cite{liu2021orthogonal,lezcano2019cheap}. Later sections elaborate on efficient approaches to uphold this orthogonality.

\subsection{Efficient Orthogonal Parameterization}\label{sect:eff_ortho}

Conventional orthogonalization procedures like the Gram-Schmidt process, although differentiable, typically entail substantial computational overhead~\cite{liu2021orthogonal}. To enhance efficiency, we utilize the Cayley transform as our parameterization technique for producing orthogonal matrices. This involves defining the orthogonal matrix as $\thickmuskip=2mu \medmuskip=2mu \bm{R} = (\bm{I} + \bm{Q})(I - \bm{Q})^{-1}$, where $\bm{Q}$ is constructed as a skew-symmetric matrix satisfying $\thickmuskip=2mu \medmuskip=2mu \bm{Q} = -\bm{Q}^\top$. The Cayley parameterization incurs a minor limitation: it exclusively yields orthogonal matrices with determinant $1$, which are members of the special orthogonal group. However, in practice, this restriction has no negative effect on model performance. Despite the application of the Cayley transform for orthogonal matrix parameterization, the representation can still be highly inefficient in terms of parameters, especially when $d$ is large. To mitigate this inefficiency, we suggest modeling $\bm{R}$ as a block-diagonal matrix comprising $r$ blocks, thereby obtaining the structure:

\begin{equation}
    \footnotesize
    \bm{R} = \text{diag}(\bm{R}_1, \bm{R}_2, \cdots, \bm{R}_r) = \begin{bmatrix}
    \bm{R}_1 \in \text{O}(\frac{d}{r}) &  & \\
    & \ddots & \\
    & & \bm{R}_r \in \text{O}(\frac{d}{r})
    \end{bmatrix} \in \text{O}(d)
\end{equation}

Here, $\text{O}(d)$ refers to the group of orthogonal matrices in $d$-dimensional space, where $\thickmuskip=2mu \medmuskip=2mu \bm{R} \in \mathbb{R}^{d \times d}$ and, for all $i$, $\thickmuskip=2mu \medmuskip=2mu \bm{R}_i \in \mathbb{R}^{d/r \times d/r}$ are orthogonal matrices. In the special case when $\thickmuskip=2mu \medmuskip=2mu r=1$, the block-diagonal structure collapses to a standard, unconstrained orthogonal matrix. For a general orthogonal matrix of shape $\thickmuskip=2mu \medmuskip=2mu d\times d$, the parameter count is $\thickmuskip=2mu \medmuskip=2mu d(d-1)/2$, implying a computational complexity of $\mathcal{O}(d^2)$. If the orthogonal matrix is constructed as an $r$-block diagonal matrix, its number of free parameters reduces to $\thickmuskip=2mu \medmuskip=2mu d(d/r-1)/2$, so its complexity is $\thickmuskip=2mu \medmuskip=2mu \mathcal{O}(d^2/r)$. Moreover, sharing the same block across all $i$ (\emph{i.e.}, setting $\thickmuskip=2mu \medmuskip=2mu \bm{R}_i = \bm{R}_j$ for any $i \neq j$) can further shrink the number of parameters, yielding a complexity of $\mathcal{O}(d^2/r^2)$. Notably, regardless of these parameter reduction strategies, the resulting matrix $\bm{R}$ retains its orthogonality property, and thus, the hyperspherical energy is still preserved.

Let us analyze parameter efficiency of OFT in comparison to LoRA. For LoRA, which uses a low-rank parameterization with rank $r'$, the total number of trainable weights is $\thickmuskip=2mu \medmuskip=2mu r'(d+n)$. If we suppose that both $r$ and $r'$ are linearly related to the model's neuron dimensionality $d$ (for instance, $\thickmuskip=2mu \medmuskip=2mu r=r'=\alpha d$, where $\thickmuskip=2mu \medmuskip=2mu 0 < \alpha \leq 1$ is a fixed constant), LoRA’s parameter complexity becomes $\mathcal{O}(d^2+dn)$ while that for OFT is only $\mathcal{O}(d)$. For concreteness, consider a scenario where the weight matrix has dimensions $\thickmuskip=2mu \medmuskip=2mu 128 \times 128$: LoRA requires $2,048$ tunable parameters if $\thickmuskip=2mu \medmuskip=2mu r' = 8$, whereas OFT requires just $960$ parameters under $\thickmuskip=2mu \medmuskip=2mu r = 8$ with no block sharing.

\subsection{Constrained Orthogonal Finetuning}

The capacity of original OFT can be further restricted by ensuring that the updated model remains in close proximity to its pretrained counterpart. In particular, COFT operates with the following forward computation:

\begin{equation}\label{eq:coft_general}
    \footnotesize
    \bm{z} = \bm{W}^\top \bm{x} = (\bm{R} \cdot \bm{W}^0)^\top \bm{x},~~~\text{s.t.}~\bm{R}^\top \bm{R} = \bm{R}\bm{R}^\top = \bm{I},~\norm{\bm{R} - \bm{I}} \leq \epsilon
\end{equation}

which imposes both an orthogonality constraint and a bound on how much $\bm{R}$ can deviate from the identity matrix (as measured by $\epsilon$). The orthogonality requirement may be satisfied through Cayley parameterization, described in Section~\ref{sect:eff_ortho}. Imposing the additional $\epsilon$-deviation restriction directly within the Cayley parameterization framework is, however, challenging. To better understand the behavior of the Cayley transform, we

As a result of the series converging in the operator norm, the constraint $\thickmuskip=2mu \medmuskip=2mu\|\bm{R}-\bm{I}\|\leq\epsilon$ can be equivalently reframed within the Cayley transform. Specifically, it translates into an alternative constraint $\thickmuskip=2mu \medmuskip=2mu \|\bm{Q}-\bm{0}\|\leq\epsilon'$, where $\bm{0}$ represents a matrix of all zeros and $\epsilon'$ serves as a distinct hyperparameter governing the permissible error (separate from $\epsilon$). The imposed restriction on the matrix $\bm{Q}$ is straightforward to enforce via projected gradient descent. To ensure that the orthogonal matrix $\bm{R}$ is initialized as the identity, we begin with $\bm{Q}$ set to an all-zero configuration. COFT thus introduces a framework where two constraints are enforced explicitly: minimization of the change in hyperspherical energy, and limiting the extent of departure from the pretrained model. While the latter constraint is usually incorporated implicitly in conventional finetuning approaches, COFT brings it to the forefront as an explicit objective. Although OFT already achieves strong performance, our empirical results suggest that COFT further enhances the stability of finetuning due to this clear deviation constraint. As depicted in Figure~\ref{fig:eps_coft}, the parameter $\epsilon$ in COFT modulates the trade-off between flexibility and adherence to the original model. Larger values of $\epsilon$ allow COFT to closely match the behavior of OFT-finetuned models, whereas smaller $\epsilon$ make COFT-finetuned models approach the behavior of the underlying pretrained text-to-image diffusion model.

\subsection{Re-scaled Orthogonal Finetuning}

Building upon the basic OFT framework, we introduce an enhancement that allows for learning a separate scaling factor for each neuron. The rationale behind this is that scaling neuron activations leaves hyperspherical energy unchanged, since normalization removes the effect of the magnitude. Formally, our forward computation is given by $\thickmuskip=2mu \medmuskip=2mu\bm{z}=(\bm{R}\bm{W}^0\bm{D})^\top\bm{x}$\footnote{\textbf{Errata}: The NeurIPS camera ready version contains an error, listing the forward pass as $\thickmuskip=2mu \medmuskip=2mu\bm{z}=(\bm{D}\bm{R}\bm{W}^0)^\top\bm{x}$. The actual implementation, however, follows the correct version shown above, ensuring the validity of results in Appendix~\ref{app:rescaled}.}, with $\thickmuskip=2mu \medmuskip=2mu \bm{D}=\text{diag}(s_1,\cdots,s_n)\in\mathbb{R}^{n\times n}$ being a trainable diagonal matrix in which each diagonal element $s_1,\cdots,s_n$ is strictly positive. Unlike the original OFT formulation in Eq.~\eqref{eq:oft_general}, which only permits learning the orthogonal matrix $\bm{R}$, this approach makes both $\bm{D}$ and $\bm{R}$ trainable. This re-scaled OFT extends the representational capacity of OFT with only a negligible increase in parameter count. For our core experiments, we utilize the standard OFT to highlight the standalone benefits of orthogonal transformations, but our findings indicate that the re-scaled variant often yields improved results (refer to Appendix~\ref{app:rescaled}).

\section{Intriguing Insights and Discussions}

\textbf{OFT’s architecture-agnostic nature.} OFT can, in theory, be integrated into any neural network architecture. For example, although LoRA is most frequently employed with attention matrices in Transformers~\cite{hulora2022}, we apply OFT exclusively to the attention weights during our comparative studies for fairness. OFT is similarly applicable to convolutional layers, offering unique interpretability due to its block-diagonal $\bm{R}$ structure—an advantage not present in LoRA. In cases where each block of $\bm{R}$ corresponds to one input channel, the approach parallels the effect of learning depth-wise convolution kernels~\cite{chollet2017xception}. If blocks in $\bm{R}$ are shared across channels, then the orthogonal transformation is uniformly applied to all neuron channels. We present initial experiments exploring OFT-based finetuning of convolutional layers in Appendix~\ref{app:convolution}.

\textbf{Connection to LoRA}. LoRA achieves stability of the pretrained weights by incorporating a low-rank matrix, thereby limiting substantial deviations from the original parameters. In contrast, OFT imposes an orthogonal (and thus full-rank) transformation on the pretrained weights, ensuring that the core pretrained structure remains intact. The forward computation in OFT can be expressed as $\thickmuskip=2mu \medmuskip=2mu \bm{z}=(\bm{R}\bm{W}^0)^\top\bm{x}=(\bm{W}^0+(\bm{R}-\bm{I})\bm{W}^0)^\top\bm{x}$, where $\thickmuskip=2mu \medmuskip=2mu (\bm{R}-\bm{I})\bm{W}^0$ plays a similar role to the low-rank update in LoRA. Notably, because $\bm{W}^0$ generally possesses full rank, the update by OFT also yields a low-rank perturbation whenever $\thickmuskip=2mu \medmuskip=2mu \bm{R}-\bm{I}$ is low-rank. OFT introduces a diagonal block size $r$ (analogous to LoRA’s rank hyperparameter $r'$) that can reduce the parameter count during training. Fundamentally, LoRA and OFT each leverage different parameter efficiency strategies: LoRA utilizes low-rank constraints to economize parameters, while OFT achieves this via block-diagonal sparsity through orthogonal transformations.

\textbf{Why OFT converges faster}. As illustrated in Figure~\ref{fig:toy_ae}, updating neuron orientations most efficiently changes the model's semantic representation—precisely the capability that OFT facilitates. Furthermore, OFT can be interpreted as adjusting neurons on a smooth hypersphere, which improves the optimization landscape. This interpretation and the resulting empirical advantages are corroborated in \cite{liu2021orthogonal}.

\textbf{Why not minimize hyperspherical energy}. Our method departs from \cite{liu2021orthogonal} in that we do not target minimizing hyperspherical energy. For classification tasks, redundancy-free neurons are advantageous. While achieving minimum hyperspherical energy ensures a uniform arrangement of neurons on the hypersphere, such an objective is counterproductive during finetuning, as it could undermine the pretrained knowledge encoded in the weights.

\textbf{Balancing flexibility and regularization in finetuning}. Our investigation reveals a fundamental tension between flexibility and regularity in finetuning approaches. Standard finetuning, though highly flexible, suffers from limited stability and a tendency toward model collapse. In contrast, OFT, despite its simplicity, achieves a favorable trade-off between these two aspects by maintaining the angular relationships between neuron pairs. The block-diagonal formulation can be interpreted as imposing a stricter form of regularization on the orthogonal matrix.

\textbf{Inference efficiency remains unchanged}. In contrast to methods such as ControlNet, the OFT framework incurs no extra computational cost during inference for the finetuned model. During inference, one can directly combine the learned orthogonal matrix $\bm{R}$ with the pretrained weights $\bm{W}^0$ via matrix multiplication, obtaining an updated weight matrix $\thickmuskip=2mu \medmuskip=2mu \bm{W}=\bm{R}\bm{W}^0$. As a result, inference runs at the same speed as with the original pretrained network.

\section{Experiments and Results}

\textbf{Experimental configuration}. For our experiments, Stable Diffusion v1.5~\cite{rombach2022high} is used as the pretrained model for text-to-image tasks. We ensure fairness by randomly sampling generated outputs from each method. For experiments focused on subject-driven generation, we adhere primarily to the DreamBooth~\cite{ruiz2023dreambooth} methodology. For controllable image synthesis, we reference protocols from ControlNet~\cite{zhang2023adding} and T2I-Adapter~\cite{mou2023t2i}. To maintain comparability with LoRA, OFT or COFT is applied only to those layers utilized by LoRA. Additional details and experimental outcomes are available in Appendix~\ref{app:settings}.

\subsection{Subject-driven Generation}

\textbf{Experimental setup}. DreamBooth~\cite{ruiz2023dreambooth} and LoRA~\cite{hulora2022} serve as our primary baselines. Each approach utilizes the same loss function defined in DreamBooth. For both DreamBooth and LoRA, we adhere to protocols specified in their respective original publications and employ their optimal hyperparameter configurations. Supplementary results can be found in Appendix~\ref{app:settings},~\ref{app:coft_oft},~\ref{app:more_qualitative}, and~\ref{app:failure}.

\textbf{Finetuning stability and convergence}. We begin by assessing the stability and convergence rate during finetuning for DreamBooth, LoRA, OFT, and COFT. The comparative outcomes are depicted in Figure~\ref{fig:teaser} and Figure~\ref{fig:db_convergence}. It is evident that COFT and OFT exhibit robust stability throughout the finetuning of the diffusion model. At the 400-iteration mark, both DreamBooth and OFT methods effectively manage to maintain control over the generation process, yet LoRA is unable to preserve the subject’s identity. By 2000 iterations, DreamBooth generates degraded images, and LoRA incorrectly synthesizes yellow fur instead of a yellow shirt. By contrast, OFT and COFT remain consistent, producing reliably controlled outputs. These findings substantiate the rapid and steady convergence properties of our OFT framework in subject-centric generation tasks. We highlight that resilience to the precise number of finetuning iterations is a significant advantage, as it eliminates the need for extensive tuning of this parameter across different subjects. For OFT and COFT, it is feasible to employ a relatively high iteration count by default. With an appropriate choice of $\epsilon$ for COFT, both the learning rate and the iteration number can be set with minimal effort.

\textbf{Quantitative comparison}. Adhering to the evaluation methodology in \cite{ruiz2023dreambooth}, we perform a quantitative analysis to assess subject fidelity (using DINO~\cite{caron2021emerging} and CLIP-I~\cite{radford2021learning}), alignment with textual prompts (CLIP-T~\cite{radford2021learning}), and output diversity (LPIPS~\cite{zhang2018unreasonable}). CLIP-I is calculated as the mean pairwise cosine similarity of CLIP features between synthesized and real images, whereas DINO operates analogously but uses ViT S/16 DINO features. CLIP-T is derived from the average cosine similarity between the text prompt embedding and the generated images’ CLIP embeddings. To quantify diversity, we compute the average LPIPS cosine similarity across images generated for the same subject and prompt. According to Table~\ref{table:dreambooth_final}, both COFT and OFT surpass DreamBooth and LoRA by a substantial margin in DINO and CLIP-I, and attain comparable or improved scores in prompt adherence and diversity. Each method is evaluated by independently finetuning the pretrained model with 30 different random initializations, thereby mitigating the influence of stochastic variation. The results demonstrate that our OFT framework not only ensures superior convergence and stability, but also achieves consistently stronger performance at convergence.

\textbf{Qualitative comparison}. To provide a clearer sense of the advantages offered by OFT, we present a selection of randomly chosen subject-driven generation samples in Figure~\ref{fig:dreambooth_final}. In order to maintain fairness, each example is produced from the identical finetuned model across all methods; as such, text prompts are not individually adapted to maximize output quality for each method. For each technique, we report outcomes from the model yielding the top validation CLIP scores. Inspection of Figure~\ref{fig:dreambooth_final} reveals that both OFT and COFT exhibit strong abilities in maintaining semantic consistency of the subject, while LoRA frequently struggles with subject fidelity (\emph{e.g.}, LoRA is unable to retain subject identity in the case of the bowl). In addition, OFT and COFT demonstrate superior responsiveness to textual cues compared to DreamBooth, which, although often preserving subject appearance, regularly produces outputs misaligned with the given prompt (\emph{e.g.}, see the initial row for the bowl instance). These qualitative assessments indicate that OFT offers improved balance between controllability and subject preservation relative to competing approaches. Furthermore, the number of optimization steps does not markedly affect OFT’s performance; the framework remains robust with increased iteration counts, unlike DreamBooth and LoRA, whose results may degrade. Additional qualitative visualizations are provided in Appendix~\ref{app:more_qualitative}. A user study, detailed in Appendix~\ref{app:human}, corroborates OFT’s strengths.

\subsection{Controllable Generation}

\textbf{Settings}. For baseline comparison, we adopt ControlNet~\cite{zhang2023adding}, T2I-Adapter~\cite{mou2023t2i}, and LoRA~\cite{hulora2022}. The main experiments encompass three demanding conditional generation scenarios: Canny edge to image~(C2I) on COCO~\cite{lin2014microsoft}, segmentation map to image~(S2I) using ADE20K~\cite{zhou2017scene}, and landmark to face~(L2F) leveraging CelebA-HQ~\cite{karras2018progressive,xia2021tedigan}. All approaches are employed to finetune Stable Diffusion~(SD) v1.5 on the respective datasets for 20 training epochs. Supplementary experimental findings can be found in Appendix~\ref{app:more_qualitative},\ref{app:more_control_task},\ref{app:failure}.

\textbf{Convergence}. We assess the rate of convergence for ControlNet, T2I-Adapter, LoRA, and COFT within the context of the L2F task, utilizing both numerical and visual analyses. As for the evaluation metric, we calculate the average $\ell_2$ distance between the predicted and control face landmarks. Figure~\ref{fig:face_convergence} displays the progression of face landmark error as each model undergoes finetuning over increasing numbers of epochs. Our findings indicate that COFT and OFT demonstrate notably accelerated convergence compared to other approaches. For instance, LoRA requires 20 epochs to attain the level of performance that OFT achieves by the 8-th epoch. It is worth highlighting that both OFT and COFT employ a comparable number of trainable parameters to LoRA (which is substantially fewer than ControlNet), yet they are able to converge considerably more efficiently than current baselines. Further supporting the rapid convergence of OFT, Figure~\ref{fig:teaser} shows that OFT surpasses ControlNet and LoRA in data efficiency, achieving strong convergence using merely 5\% of the ADE20K dataset. In our qualitative evaluation, we compare OFT, COFT, and ControlNet, as ControlNet is the closest baseline in terms of landmark accuracy. According to Figure~\ref{fig:face_convergence_qual}, both OFT and COFT display reliable convergence, with generated face poses aligning incrementally with the control landmarks. In contrast, ControlNet demonstrates abrupt improvement in controllability after epoch 8, which aligns with the phenomenon of sudden convergence described in \cite{zhang2023adding}. The consistent and predictable convergence behavior of OFT facilitates its practical application, rendering the process of selecting appropriate hyperparameters more straightforward.

\textbf{Quantitative comparison}. To assess the effectiveness of controllable generation, we propose a metric for measuring control consistency. This metric involves extracting the control signal from the output image and comparing it to the original input control. Specifically, for the C2I task, IoU and F1 score are used; for the S2I task, we report mean IoU, as well as mean and overall accuracy; for the L2F task, we measure the mean $\ell_2$ distance between predicted landmarks and the corresponding control landmarks. Further details regarding these consistency metrics can be found in Appendix~\ref{app:settings}. Each competing method was evaluated with its optimal hyperparameters. The quantitative outcomes presented in Table~\ref{table:control_gen} indicate that both OFT and COFT deliver substantially stronger and more precise control relative to other approaches. It is also observed that adapter-based methods (\emph{e.g.}, T2I-Adapter and ControlNet) tend to have slower convergence and ultimately underperform. When comparing LoRA with ControlNet, LoRA yields superior results in the S2I task but falls short on the C2I and L2F tasks. Generally, even after extensive hyperparameter optimization, LoRA’s overall performance remains capped at a relatively low level. In contrast, the OFT framework continues to improve as the parameter count increases, suggesting that its potential has not been fully realized. Importantly, our framework for quantitative assessment in controllable generation represents a significant innovation, providing reliable evaluation of control fidelity in finetuned models—subject to the accuracy limits of the employed segmentation or detection backbone.

\textbf{Qualitative comparison}. In addition to quantitative metrics, we conduct a qualitative analysis of OFT and COFT, contrasting them with leading approaches such as ControlNet, T2I-Adapter, and LoRA. As depicted by the randomly selected samples in Figure~\ref{fig:control_final}, both OFT and COFT generate images that are not only visually realistic and high in fidelity, but also exhibit precise adherence to input control conditions. Examining the S2I task, it becomes evident that LoRA is unable to reliably synthesize images based on the given segmentation map, whereas ControlNet, OFT, and COFT are capable of effective control. Unlike ControlNet, our methods produce images featuring richer detail and more plausible geometric arrangements, all while using substantially fewer model parameters. In the context of the C2I task, OFT and COFT demonstrate the capability to infer semantically coherent images from rudimentary Canny edge inputs, outperforming both T2I-Adapter and LoRA, which deliver notably inferior results. Within the L2F task, our approach yields the most precise pose-controlled face generations, even in scenarios involving complex facial poses. Across all three control scenarios, OFT and COFT consistently surpass existing state-of-the-art methods in qualitative image quality, thus substantiating the efficacy of the OFT framework for controllable image synthesis. For further in-depth qualitative insights, we include additional visual examples for each control task in Appendix~\ref{app:control_exp}. Additionally, we show in Appendix~\ref{app:more_control_task} that OFT extends effectively to other control tasks, such as dense pose to human body, sketch to image, and depth to image generation.

\section{Concluding Remarks and Open Problems}

Drawing inspiration from the critical role of angular relationships among neurons in shaping visual semantics, we introduce orthogonal finetuning as a straightforward yet powerful technique for refining text-to-image diffusion models. Our focus centers on two primary applications: subject-driven generation and controllable generation. In comparison to prior approaches, OFT achieves superior control and greater stability during finetuning, all while requiring a reduced number of adjustable parameters. Significantly, OFT avoids any additional computational overhead during inference, resulting in a model that is both effective and practical for deployment.

The adoption of OFT gives rise to several compelling open questions. Firstly, OFT enforces orthogonality via Cayley parametrization, a process that necessitates the computation of a matrix inverse. This requirement imposes some constraints on the scalability of OFT. We mitigate this to some degree using block diagonal parametrization, but finding methods to accelerate matrix inversion in a differentiable context remains an unresolved issue. Secondly, OFT exhibits distinct potential with respect to compositionality: the orthogonal matrices yielded from various OFT finetuning tasks can be sequentially multiplied and the product remains orthogonal. Whether this ensemble of orthogonal matrices can jointly retain information from multiple downstream tasks represents a promising avenue for future investigation. Lastly, OFT's parameter efficiency largely hinges on the block diagonal structure, a design choice that introduces bias and limits versatility. Developing approaches to enhance parameter efficiency in a less restrictive and more robust fashion continues to be an open and important research problem.

\end{document}